\section{Related work}
\textbf{Design-based estimation under interference.} \citet{aronow2017} define exposure mappings and give Horvitz--Thompson (HT) estimators; \citet{ugander2013} propose graph cluster randomisation to make full-neighbourhood exposure likely, and \citet{eckles2017} study its bias--variance trade-off and analysis methods using neighbour treatment information; \citet{karrer2021} describe a deployed system. \citet{leung2022} derives rate-optimal cluster designs for HT estimation of the GATE under spatial interference, and \citet{savje2021} characterise what standard estimators estimate under unknown interference. \citet{basse2018} show limits of design-based inference under arbitrary interference. \citet{cortez2022} give staggered-rollout designs without network knowledge, and \citet{yu2022} estimate the total treatment effect with unknown network structure.\\
\textbf{Regression-based estimation.} \citet{chin2019} proposes regression adjustment that predicts all-treated and all-control outcomes, with properties under linear models; \citet{han2023} select graph-based and treatment-based covariates and use block bootstrap; \citet{liwager2022} analyses estimators under random-graph asymptotics. \citet{baird2018} study randomised saturation designs, a related design idea we include. \citet{covdesign2023} optimise the assignment covariance for HT estimators in a one-hop model.\\
\textbf{Positioning.} Prior regression work establishes validity under a correctly specified model or provides methodology. We instead ask how large the error is, and \emph{what it is made of}, across mechanisms and graphs, give an exact weight representation with post-assignment diagnostics, and map when linear and flexible regressions fail (extrapolation, high degree, poor clusters, wrong graph). Graph neural networks \citep{kipf2017,defferrard2016} are an alternative for learning exposure features; we do not evaluate them (\S\ref{sec:disc}).
